% =====================================================================
%  Sezione 2.5 — Le tre rappresentazioni geometriche
%  Sostituisce le sottosezioni "Metodo A / A2 / A3" della versione
%  precedente. Il Metodo B resta in una sezione separata.
% =====================================================================

\section{Tre rappresentazioni geometriche dello stesso movimento}
\label{sec:rappresentazioni}

Dai trentatré \emph{landmark} restituiti dallo stimatore di posa si possono
costruire descrizioni molto diverse dello stesso movimento. In questo lavoro ne
sono state considerate tre, che non sono tre modelli in competizione ma tre
punti di vista sullo stesso dato: insieme formano il bacino di candidate da cui
la selezione descritta nella Sezione~\ref{sec:dtw} estrae il sottoinsieme
effettivamente utilizzato.

\begin{itemize}
    \item \textbf{A1}, la sagoma come ellisse: il corpo trattato come nuvola di
          punti nel piano immagine, senza conoscenza dell'anatomia;
    \item \textbf{A2}, la geometria in tre dimensioni: l'unica delle tre che non
          sia una proiezione, costruita sui \emph{world landmark};
    \item \textbf{A3}, la geometria riferita al corpo: tutto normalizzato
          rispetto alla lunghezza del tronco, quindi adimensionale.
\end{itemize}

I nomi delle grandezze sono prefissati con quello della rappresentazione di
appartenenza. Non è una precauzione formale: nella versione iniziale del lavoro
la grandezza \texttt{d\_trunk\_angle} esisteva sia in A2 sia in A3 con unità di
misura diverse — gradi per fotogramma contro gradi al secondo, con mediane
$0{,}167$ e $1{,}700$ — ed era l'unica colonna in comune fra i due insiemi. Un
accorpamento senza prefisso le avrebbe fuse silenziosamente.

\subsection{A1 --- la sagoma come ellisse}

Sette \emph{landmark} — naso, spalle, anche e caviglie — vengono trattati come
una distribuzione di punti nel piano immagine. Dalla matrice di covarianza si
ricavano autovalori e autovettori: gli autovettori danno le direzioni principali
lungo cui il corpo si distribuisce, gli autovalori la dispersione lungo tali
direzioni.

Intuitivamente il corpo viene approssimato da un'ellisse: l'asse maggiore
fornisce l'orientamento dominante, il rapporto fra gli assi descrive quanto la
configurazione appaia allungata o compressa. Le grandezze estratte sono i due
assi, il loro rapporto, l'inclinazione dell'asse maggiore, le rispettive
derivate e le due componenti della velocità del baricentro della nuvola.

Il punto di forza di questa rappresentazione è che \emph{non conosce
l'anatomia}: non le serve identificare correttamente il tronco per accorgersi
che la sagoma è passata da verticale e stretta a orizzontale e larga. Funziona
quindi anche quando lo stimatore sbaglia l'attribuzione dei singoli giunti,
purché la nuvola complessiva sia nella posizione giusta. Il limite è lo stesso:
non distingue una persona distesa da una chinata in avanti, perché entrambe
producono un'ellisse orizzontale.

Una precisazione implementativa. L'inclinazione dell'asse di un'ellisse è una
\emph{direzione}, non un verso: ha periodo $180^\circ$ e non $360^\circ$.
Trattarla con il periodo sbagliato produce salti artificiali nella serie
temporale — è stato misurato un passaggio da $201{,}9^\circ$ a $49{,}5^\circ$
fra due fotogrammi consecutivi, cioè una rotazione fisicamente impossibile — che
si propagano alla derivata e a ogni grandezza accumulata.

\subsection{A2 --- la geometria in tre dimensioni}

La seconda rappresentazione utilizza i \emph{world landmark}, cioè le coordinate
metriche tridimensionali stimate dalla rete. È l'unica delle tre a non essere
una proiezione, e questo la rende potenzialmente meno sensibile al punto di
vista della telecamera.

Dal punto medio delle spalle e da quello delle anche si ricava il vettore del
tronco, dal quale si calcola l'angolo rispetto alla verticale del riferimento;
il naso fornisce la posizione della testa. Le grandezze sono l'angolo del
tronco, la sua derivata con segno e in modulo, la velocità della testa e
l'altezza della testa sopra il bacino, normalizzata sulla lunghezza del tronco.

Due precisazioni sono necessarie, e riguardano entrambe il sistema di
riferimento.

La prima è che il riferimento dei \emph{world landmark} è \emph{canonico e
centrato sulle anche}, non allineato alla gravità reale. L'angolo così calcolato
non è quindi un'inclinazione assoluta rispetto al pavimento: è una misura
relativa al corpo. Resta informativa sui dati, ma chiamarla «inclinazione
rispetto alla verticale» sarebbe impreciso.

La seconda è più stringente. Poiché il riferimento è centrato sulle anche, il
bacino \textbf{è l'origine e non si muove per costruzione}. La misura lo
conferma: su quaranta fotogrammi consecutivi la deviazione standard della sua
posizione è di $3\cdot10^{-5}$ metri, cioè tre centesimi di millimetro, contro i
$1{,}75\cdot10^{-2}$ metri del naso — quattrocento volte tanto. Le grandezze
relative alla traslazione del bacino, presenti nella formulazione iniziale,
misuravano quindi rumore numerico e sono state rimosse. La traslazione assoluta
del corpo con questa rappresentazione non è misurabile: la misurano A1 e A3, nel
piano immagine.

\subsection{A3 --- la geometria riferita al corpo}

La terza rappresentazione utilizza i \emph{landmark} bidimensionali come A1, ma
con una differenza decisiva: ogni grandezza viene divisa per la lunghezza del
tronco. Il risultato è adimensionale, quindi indipendente dalla distanza della
persona dalla telecamera e dalla sua corporatura: la stessa posa a due o a
cinque metri, su una persona alta o bassa, produce gli stessi numeri.

Il centro del bacino è l'origine del riferimento locale. Le grandezze sono la
lunghezza del tronco, l'angolo del tronco nel piano immagine e la sua derivata,
le altezze di testa, ginocchia e caviglie rispetto al bacino, la larghezza delle
spalle, la velocità del bacino in lunghezze di torso al secondo e la velocità
della testa relativa al bacino.

È la rappresentazione anatomicamente più ricca — l'unica che osservi gli arti
inferiori, e sapere che le caviglie sono salite all'altezza del bacino
distingue una caduta all'indietro da un semplice piegarsi — ma vive nel piano
immagine, e ne eredita quindi la dipendenza dal punto di vista.

Va inoltre osservato che la normalizzazione rispetto al corpo, che è il pregio
di questa rappresentazione, ne è anche il limite: normalizzando tutto rispetto
al bacino, A3 \emph{cancella per costruzione} l'informazione su dove il corpo si
trovi nella stanza. È il motivo per cui viene affiancata dalla grandezza
descritta nella Sezione~\ref{sec:pavimento}.

\subsection{Cosa è stato rimosso, e perché}

Rispetto alla formulazione iniziale, che prevedeva dieci grandezze per ciascuna
rappresentazione, alcune sono state rimosse. Le ragioni sono di tre tipi.

\textbf{Grandezze degenerate per costruzione.} Le due relative alla traslazione
del bacino in A2, per la ragione discussa sopra.

\textbf{Grandezze dimensionalmente prive di senso.} Ciascuna rappresentazione
conteneva una grandezza denominata \emph{energia}, calcolata come somma di
quadrati dei contributi dinamici. Quelle somme sommavano quadrati di gradi, di
rapporti adimensionali, di pixel e di frazioni di immagine: quantità non
omogenee, il cui risultato non ha un'unità interpretabile. Sul dataset
presentavano rapporti fra valore massimo e mediano dell'ordine di $7\cdot10^5$,
ed erano all'origine della disparità di scala fra le colonne in ingresso — un
fattore 646 fra la più grande e la più piccola — che senza normalizzazione
rendeva alcune grandezze invisibili al modello.

\textbf{Derivate seconde.} Le accelerazioni presenti in A2 sono state rimosse
per la ragione descritta nella Sezione~\ref{sec:rumore}: se la derivata prima
amplifica il rumore dello stimatore di posa, la derivata seconda lo amplifica
due volte.

Il bacino risultante conta trentadue grandezze: nove per A1, cinque per A2, nove
per A3 e la terna integrale di ciascuna rappresentazione, descritta nella
Sezione~\ref{sec:integrali}.

\subsection{Una convenzione temporale unica}

Un'ultima precisazione riguarda le derivate. Nella formulazione iniziale A1
calcolava differenze \emph{per fotogramma} mentre A2 e A3 dividevano per
$\Delta t$: confrontare A1 con A3 significava quindi confrontare anche due
convenzioni temporali diverse, e cambiando frame rate una sola delle due
cambiava scala. Nell'implementazione attuale tutte le derivate sono espresse per
secondo.
